\documentclass[11pt]{article}
\usepackage[margin=1in]{geometry}
\usepackage{amsmath, amssymb, amsthm}
\usepackage{algorithm}
\usepackage{algorithmic}
\usepackage{bm}

\title{Efficient Alternating Optimization for RKHS-Constrained CP Decomposition with Missing Data}
\author{}
\date{}

\begin{document}

\maketitle

\section{Problem Formulation}

We consider the mode-$k$ subproblem in the Alternating Least Squares (ALS) optimization of a rank-$r$ CP decomposition for a $d$-way tensor $\mathcal{T} \in \mathbb{R}^{n_1 \times \dots \times n_d}$. The tensor contains missing entries, and the set of observed indices is denoted by $\Omega$, with $|\Omega| = q$. The total number of entries is $N = \prod n_i$. We assume $q \ll N$.

The mode-$k$ factor matrix $A_k \in \mathbb{R}^{n \times r}$ (where $n = n_k$) is constrained to lie in a Reproducing Kernel Hilbert Space (RKHS) generated by a positive semidefinite kernel matrix $K \in \mathbb{R}^{n \times n}$. We parameterize the factor as $A_k = KW$, where $W \in \mathbb{R}^{n \times r}$ is the unknown weight matrix. 

Given fixed factor matrices for all other modes, we define $Z \in \mathbb{R}^{M \times r}$ as their Khatri-Rao product ($M = \prod_{i \ne k} n_i$). The regularized least squares problem for $W$ yields the following linear system of size $nr \times nr$:
\begin{equation}
    \mathcal{A} \vecop(W) = \mathbf{b},
    \label{eq:system}
\end{equation}
where the system matrix $\mathcal{A}$ and the right-hand side $\mathbf{b}$ are defined as:
\begin{align}
    \mathcal{A} &= (Z \otimes K)^T S S^T (Z \otimes K) + \lambda (I_r \otimes K), \\
    \mathbf{b} &= (I_r \otimes K) \vecop(B).
\end{align}
Here, $S \in \mathbb{R}^{N \times q}$ is the selection matrix for the observed entries, and $B = T Z$ is the Matricized Tensor Times Khatri-Rao Product (MTTKRP) of the mode-$k$ unfolding $T$ (with missing entries set to zero). 

Direct solution of this system requires $O(n^3 r^3)$ operations, which is computationally prohibitive. Furthermore, explicitly forming the matrix $\mathcal{A}$ or the dense factor $Z$ involves $O(N)$ memory/time, which is infeasible. We propose an efficient solution using the Preconditioned Conjugate Gradient (PCG) method.

\section{Method: Preconditioned Conjugate Gradient}

The PCG method is an iterative solver for symmetric positive definite systems. Its efficiency in this context relies on two components: performing the matrix-vector product (MVP) in time proportional to $q$ rather than $N$, and using a preconditioner that clusters the eigenvalues of the system matrix.

\subsection{Efficient Matrix-Vector Multiplication}
\label{sec:mvp}

We need to compute $\mathbf{y} = \mathcal{A} \mathbf{x}$ for an arbitrary vector $\mathbf{x} \in \mathbb{R}^{nr}$. Let $\mathbf{x} = \vecop(X)$ for $X \in \mathbb{R}^{n \times r}$. We decompose the operation into a regularization term and a data fidelity term.

\subsubsection{Regularization Term}
The term $\mathbf{y}_{\text{reg}} = \lambda (I_r \otimes K) \vecop(X)$ simplifies using the Kronecker product identity $(B^T \otimes A) \vecop(C) = \vecop(ACB)$:
\begin{equation}
    \mathbf{y}_{\text{reg}} = \lambda \vecop(K X).
\end{equation}
Computationally, we form the matrix product $KX$.
\begin{itemize}
    \item \textbf{Complexity:} $O(n^2 r)$.
\end{itemize}

\subsubsection{Data Fidelity Term}
The term is $\mathbf{y}_{\text{data}} = (Z \otimes K)^T S S^T (Z \otimes K) \vecop(X)$. Using the identity $(Z \otimes K) \vecop(X) = \vecop(K X Z^T)$, this term can be interpreted as projecting the model tensor onto the observed set $\Omega$, applying the mask, and projecting back. To avoid $O(N)$ costs, we perform these operations sparsely.

\begin{enumerate}
    \item \textbf{Map to Factor Space:} Compute $U = K X \in \mathbb{R}^{n \times r}$. 
    Cost: $O(n^2 r)$.
    
    \item \textbf{Sparse Evaluation (Forward Projection):} The dense tensor model is $\mathcal{M} = U Z^T$. We only need its values at indices in $\Omega$. For each observed index $m \in \{1, \dots, q\}$, let $(i_m, j_m)$ be the corresponding indices in the mode-$k$ unfolding (where $i_m \in \{1,\dots,n\}$ and $j_m \in \{1,\dots,M\}$).
    The value is:
    \begin{equation}
        v_m = (U)_{i_m, :} (Z)_{j_m, :}^T = \sum_{\ell=1}^r U_{i_m, \ell} Z_{j_m, \ell}.
    \end{equation}
    The row $Z_{j_m, :}$ is the element-wise product of the rows of the fixed factor matrices corresponding to the multi-index $j_m$. It can be computed on-the-fly in $O(dr)$ operations.
    Total Cost: $O(qdr)$.
    
    \item \textbf{Sparse Adjoint (Backward Projection):} We compute $\mathbf{u} = (Z \otimes K)^T S S^T \vecop(\mathcal{M})$. Let $\mathbf{v} \in \mathbb{R}^q$ be the vector of values computed above. This step is equivalent to computing $\vecop(K H Z)$, where $H \in \mathbb{R}^{n \times M}$ is a sparse matrix with entries $v_m$ at $(i_m, j_m)$.
    First, compute $Y = H Z \in \mathbb{R}^{n \times r}$. Since $H$ is sparse, we iterate over the $q$ non-zeros:
    \begin{equation}
        Y_{i, :} \leftarrow Y_{i, :} + v_m Z_{j_m, :}, \quad \forall m=1,\dots,q.
    \end{equation}
    Again, $Z_{j_m, :}$ is generated on-the-fly.
    Total Cost: $O(qdr)$.
    
    \item \textbf{Kernel Smoothing:} Finally, compute $Y_{\text{data}} = K Y$. 
    Cost: $O(n^2 r)$.
\end{enumerate}

The final result is $\mathbf{y} = \vecop(Y_{\text{data}} + \lambda U)$. The total complexity per MVP is $O(qdr + n^2 r)$.

\subsection{Preconditioner}

The convergence of PCG depends on the spectral properties of $\mathcal{A}$. The projection $S S^T$ disrupts the Kronecker structure. We use a standard approximation for random sampling, $S S^T \approx \rho I_N$, where $\rho = q/N$ is the sampling ratio.
Substituting this into $\mathcal{A}$, we obtain the preconditioner $\mathcal{P}$:
\begin{align}
    \mathcal{P} &= \rho (Z \otimes K)^T (Z \otimes K) + \lambda (I_r \otimes K) \nonumber \\
    &= \rho (Z^T Z \otimes K^2) + \lambda (I_r \otimes K).
\end{align}
Here we utilized the property $(A \otimes B)^T (A \otimes B) = A^T A \otimes B^T B$ and the symmetry of $K$. Let $G = Z^T Z \in \mathbb{R}^{r \times r}$. $G$ is efficiently computed as the Hadamard product of the Gram matrices of the other modes: $G = \mathop{\circledast}_{m \neq k} (A_m^T A_m)$.
Factoring out $(I_r \otimes K)$ on the left:
\begin{equation}
    \mathcal{P} = (I_r \otimes K) \left[ \rho (G \otimes K) + \lambda I_{nr} \right].
\end{equation}

\subsubsection{Applying the Preconditioner}
We need to solve $\mathcal{P} \mathbf{z} = \mathbf{r}$ for $\mathbf{z}$. This requires inverting $\mathcal{P}$. We exploit the eigendecompositions of the small matrices $K$ and $G$.
Let $K = Q_K \Lambda_K Q_K^T$ and $G = Q_G \Lambda_G Q_G^T$, where $\Lambda_K = \text{diag}(\mu_1, \dots, \mu_n)$ and $\Lambda_G = \text{diag}(\sigma_1, \dots, \sigma_r)$.
The matrix $\mathcal{P}$ shares the eigenvector basis $Q_G \otimes Q_K$. The eigenvalues of $\mathcal{P}$ are:
\begin{equation}
    \xi_{ij} = \rho \sigma_j \mu_i^2 + \lambda \mu_i, \quad i=1\dots n, \, j=1\dots r.
\end{equation}
The application of $\mathcal{P}^{-1}$ to a vector $\mathbf{r} = \vecop(R)$ proceeds as follows:
\begin{enumerate}
    \item \textbf{Basis Rotation:} Compute $\tilde{R} = Q_K^T R Q_G$. 
    Cost: $O(n^2 r + nr^2)$.
    \item \textbf{Spectral Scaling:} Compute $\tilde{Z}_{ij} = \tilde{R}_{ij} / \xi_{ij}$. If $\xi_{ij}$ is close to zero (e.g., if $K$ is rank-deficient), we use the pseudo-inverse (set $\tilde{Z}_{ij}=0$).
    Cost: $O(nr)$.
    \item \textbf{Inverse Rotation:} Compute $Z_{\text{mat}} = Q_K \tilde{Z} Q_G^T$. Then $\mathbf{z} = \vecop(Z_{\text{mat}})$.
    Cost: $O(n^2 r + nr^2)$.
\end{enumerate}

\section{Complexity Analysis}

The algorithm consists of a setup phase and an iterative phase. We assume $n, r < q \ll N$.

\paragraph{Setup Costs:}
\begin{itemize}
    \item Computing $G$: $O(d r^2)$.
    \item Eigendecomposition of $K$ ($n \times n$) and $G$ ($r \times r$): $O(n^3 + r^3)$.
    \item Computing RHS $\mathbf{b}$: This involves a sparse MTTKRP similar to the MVP data term. Cost: $O(qdr + n^2 r)$.
\end{itemize}
Total Setup: $O(n^3 + qdr)$.

\paragraph{Per-Iteration Costs:}
\begin{itemize}
    \item Matrix-Vector Product: $O(qdr + n^2 r)$.
    \item Preconditioner Application: $O(n^2 r + nr^2)$.
    \item CG Vector Operations: $O(nr)$.
\end{itemize}
Total Per-Iteration: $O(qdr + n^2 r)$.

This complexity is linear in the number of observations $q$ and cubic in the mode dimension $n$ (only for setup), completely avoiding any dependency on the full tensor size $N$. This makes the method highly efficient for large, sparse tensors.

\begin{algorithm}[H]
\caption{PCG for RKHS-Constrained Mode-$k$ Update}
\begin{algorithmic}[1]
\STATE \textbf{Input:} $\Omega, \text{values}, K, \{A_m\}_{m \neq k}, \lambda$.
\STATE \textbf{Preprocessing:}
\STATE \quad $K = Q_K \Lambda_K Q_K^T$.
\STATE \quad $G = \mathop{\circledast}_{m \neq k} (A_m^T A_m) = Q_G \Lambda_G Q_G^T$.
\STATE \quad Compute $\mathbf{b} = \vecop(K \cdot \text{SparseMTTKRP}(T, Z))$.
\STATE \textbf{Initialize:} $W^{(0)}$, $\mathbf{r}^{(0)} = \mathbf{b} - \mathcal{A}\vecop(W^{(0)})$, $\mathbf{p}^{(0)} = \mathcal{P}^{-1} \mathbf{r}^{(0)}$.
\WHILE{not converged}
    \STATE $\mathbf{v} = \mathcal{A} \mathbf{p}$ \COMMENT{Using sparse MVP (Sec. \ref{sec:mvp})}
    \STATE $\alpha = (\mathbf{r}^T \mathcal{P}^{-1} \mathbf{r}) / (\mathbf{p}^T \mathbf{v})$
    \STATE $\vecop(W) \leftarrow \vecop(W) + \alpha \mathbf{p}$
    \STATE $\mathbf{r}_{\text{new}} = \mathbf{r} - \alpha \mathbf{v}$
    \STATE $\beta = (\mathbf{r}_{\text{new}}^T \mathcal{P}^{-1} \mathbf{r}_{\text{new}}) / (\mathbf{r}^T \mathcal{P}^{-1} \mathbf{r})$
    \STATE $\mathbf{p} \leftarrow \mathcal{P}^{-1} \mathbf{r}_{\text{new}} + \beta \mathbf{p}$
    \STATE $\mathbf{r} \leftarrow \mathbf{r}_{\text{new}}$
\ENDWHILE
\STATE \textbf{Output:} $W$.
\end{algorithmic}
\end{algorithm}

\end{document}
